\documentclass[11pt,a4paper]{article}
\usepackage[T1]{fontenc}
\usepackage[utf8]{inputenc}
\usepackage[spanish]{babel}
\usepackage{lmodern}
\usepackage[margin=2.5cm]{geometry}
\usepackage{amsmath,amssymb,amsthm}
\usepackage{titlesec}
\usepackage{needspace}
\setlength{\parindent}{0pt}
\setlength{\parskip}{5pt}
\raggedbottom
\clubpenalty=10000
\widowpenalty=10000
\titleformat{\section}{\large\bfseries}{}{0pt}{}
\titlespacing*{\section}{0pt}{14pt}{7pt}
\begin{document}
\begin{center}
{\LARGE\bfseries Acciones de grupos y conteo}\\[6pt]
{\large Ejercicio 1: incisos (a) y (b)}
\end{center}
\medskip
\section*{(a) Un coeficiente binomial no divisible por $p$}

\textbf{Enunciado.} Sean $p$ primo, $k\geq1$ y $n\geq0$, con $\gcd(p,k)=1$. Demostrar que
\[
p\nmid\binom{p^nk}{p^n}.
\]

\textbf{Demostración.} Si $n=0$, el coeficiente binomial es $\binom{k}{1}=k$, y el resultado es inmediato. Supongamos $n\geq1$.

\section*{1. Elección del grupo y del conjunto}
Tomamos
\[
G=\mathbb Z/p^n\mathbb Z,\qquad X=G\times\{1,\ldots,k\}.
\]
Un elemento de $X$ es un par $(x,i)$, con $x\in G$ y $1\leq i\leq k$. Por lo tanto,
\[
|G|=p^n,\qquad |X|=|G|k=p^nk.
\]
Podemos elegir este conjunto particular porque la cantidad de subconjuntos de un tamaño dado depende solamente del cardinal del conjunto que los contiene.

Definimos
\[
\Omega=\{S\subseteq X:|S|=p^n\}.
\]
Es decir, cada elemento $S$ de $\Omega$ es un subconjunto de $X$ con $p^n$ elementos. Por la interpretación combinatoria del coeficiente binomial,
\[
\boxed{|\Omega|=\binom{p^nk}{p^n}}.
\]

\section*{2. La acción sobre los subconjuntos}
Consideramos la acción de $G$ sobre $X$ dada por
\[
g\cdot(x,i)=(g+x,i).
\]
La suma tiene sentido porque $g$ y $x$ pertenecen al mismo grupo $G$, cuya operación es la suma módulo $p^n$.

Sobre $\Omega$ consideramos la acción inducida
\[
g\cdot S=\{(g+x,i):(x,i)\in S\}.
\]
Ya sabemos que esta regla define una acción de $G$ sobre $\Omega$.

\newpage
\section*{3. Los puntos fijos que queremos determinar}
Por definición,
\[
\Omega^G:=\{S\in\Omega:g\cdot S=S\text{ para todo }g\in G\}.
\]
Aquí un punto fijo es un \emph{subconjunto} $S$ que permanece igual como conjunto bajo la acción de todo $g\in G$. Sus elementos individuales pueden cambiar de lugar.

Para cada $i\in\{1,\ldots,k\}$, escribimos
\[
B_i:=G\times\{i\}.
\]
La igualdad que debemos demostrar es
\[
\Omega^G=\{B_1,\ldots,B_k\}.
\]
La probamos mediante las dos inclusiones.

\section*{4. Primera inclusión: $\{B_1,\ldots,B_k\}\subseteq\Omega^G$}
Fijemos $i\in\{1,\ldots,k\}$. Como hay exactamente un par $(x,i)$ por cada $x\in G$,
\[
|B_i|=|G\times\{i\}|=|G|\cdot1=p^n.
\]
Además, $B_i\subseteq X$, de modo que $B_i\in\Omega$.

Sea $g\in G$. Veamos que $g\cdot B_i=B_i$:
\begin{itemize}
\item Todo elemento de $g\cdot B_i$ tiene la forma $(g+x,i)$, con $x\in G$. Como $g+x\in G$, pertenece a $B_i$. Así,
\[
g\cdot B_i\subseteq B_i.
\]
\item Recíprocamente, sea $(y,i)\in B_i$. Como $y-g\in G$, tenemos $(y-g,i)\in B_i$, y entonces
\[
(y,i)=g\cdot(y-g,i)\in g\cdot B_i.
\]
Por lo tanto,
\[
B_i\subseteq g\cdot B_i.
\]
\end{itemize}
Las dos inclusiones prueban que $g\cdot B_i=B_i$. Como esto vale para todo $g\in G$, resulta $B_i\in\Omega^G$.

Dado que $i$ era arbitrario, concluimos
\[
\{B_1,\ldots,B_k\}\subseteq\Omega^G.
\]
Esto prueba que los subconjuntos indicados son fijos; todavía debemos descartar que existan otros.

\newpage
\section*{5. Segunda inclusión: $\Omega^G\subseteq\{B_1,\ldots,B_k\}$}
Sea $S\in\Omega^G$ arbitrario. Por definición,
\[
S\subseteq X,\qquad |S|=p^n,\qquad g\cdot S=S\quad\text{para todo }g\in G.
\]
Como $|S|>0$, elegimos un elemento $(x,i)\in S$. Vamos a demostrar que $S=B_i$.

Sea $y\in G$ arbitrario. Tomando $g=y-x\in G$, obtenemos
\[
(y,i)=(g+x,i)=g\cdot(x,i)\in g\cdot S=S.
\]
Como esto vale para todo $y\in G$, se sigue que
\[
B_i=G\times\{i\}\subseteq S.
\]
Ambos conjuntos tienen $p^n$ elementos:
\begin{itemize}
\item $|B_i|=|G|=p^n$, porque su segunda coordenada es fija;
\item $|S|=p^n$, porque $S\in\Omega$ y así se definió $\Omega$.
\end{itemize}
Una inclusión entre dos conjuntos finitos del mismo cardinal es una igualdad. Por consiguiente,
\[
S=B_i.
\]
Esto muestra que todo $S\in\Omega^G$ pertenece a $\{B_1,\ldots,B_k\}$, es decir,
\[
\Omega^G\subseteq\{B_1,\ldots,B_k\}.
\]
La igualdad $S=B_i$ es la que permite concluir que no hay otros subconjuntos fijos.

\section*{6. Cantidad de puntos fijos}
Las dos inclusiones anteriores prueban
\[
\Omega^G=\{B_1,\ldots,B_k\}.
\]
Estos $k$ subconjuntos son distintos: si $i\neq j$, entonces $(0,i)\in B_i$ pero $(0,i)\notin B_j$, pues la segunda coordenada de todos los elementos de $B_j$ es $j$. Por lo tanto,
\[
\boxed{|\Omega^G|=k}.
\]

\section*{7. Aplicación del lema de puntos fijos}
Por el lema de los puntos fijos,
\[
|\Omega|\equiv|\Omega^G|\pmod p.
\]
Sustituyendo los cardinales calculados,
\[
\binom{p^nk}{p^n}\equiv k\pmod p.
\]
Como $\gcd(p,k)=1$, tenemos $k\not\equiv0\pmod p$. En consecuencia,
\[
\boxed{p\nmid\binom{p^nk}{p^n}}.
\]
\hfill$\square$
\newpage
\section*{(b) Coloraciones de las caras del cubo}
\textbf{Enunciado.} Explicar, mediante acciones de grupos y la fórmula de Burnside, el conteo de coloraciones de las caras del cubo con seis colores, utilizando cada color exactamente una vez y considerando equivalentes las coloraciones que difieren por una rotación.

\textbf{Demostración.}
\section*{1. El conjunto de coloraciones}
Sean $\mathcal F$ el conjunto de las seis caras del cubo y $\mathcal C$ el conjunto de los seis colores disponibles. Una coloración que utiliza cada color exactamente una vez es una biyección $c:\mathcal F\to\mathcal C$. Sea
\[
X=\{c:\mathcal F\to\mathcal C:c\text{ es biyectiva}\}.
\]
Como $|\mathcal F|=|\mathcal C|=6$, se tiene
\[
|X|=6!.
\]

\section*{2. La acción de las rotaciones}
Sea $R$ el grupo de rotaciones del cubo. Entonces $|R|=24$ y $R$ actúa sobre $\mathcal F$. Por la acción inducida sobre funciones mediante precomposición, $R$ actúa sobre $X$ mediante
\[
(r\cdot c)(y)=c(r^{-1}\cdot y),
\qquad r\in R,\ c\in X,\ y\in\mathcal F.
\]
Las órbitas de esta acción agrupan las coloraciones que difieren por una rotación.

\section*{3. Coloraciones fijas por la identidad}
La identidad fija todas las coloraciones: para todo $c\in X$ y toda cara $y\in\mathcal F$,
\[
(\operatorname{id}\cdot c)(y)
=c(\operatorname{id}^{-1}\cdot y)=c(y).
\]
Entonces $\operatorname{id}\cdot c=c$, es decir, $c\in X^{\operatorname{id}}$. Por lo tanto,
\[
|X^{\operatorname{id}}|=|X|=6!.
\]

\newpage
\section*{4. Ninguna rotación no trivial fija una coloración}
Sea $r\in R$ una rotación no trivial. Supongamos, por contradicción, que $X^r\neq\varnothing$. Entonces existe $c\in X^r$, es decir, una coloración $c\in X$ tal que
\[
r\cdot c=c.
\]

Evaluando esta igualdad en cada cara $y\in\mathcal F$ y usando la definición de la acción, obtenemos
\[
c(y)=(r\cdot c)(y)=c(r^{-1}\cdot y).
\]

Como $c$ es biyectiva, en particular es inyectiva, y por lo tanto
\[
y=r^{-1}\cdot y.
\]
Por las propiedades de la acción,
\[
r\cdot y=r\cdot(r^{-1}\cdot y)
=(rr^{-1})\cdot y=\operatorname{id}\cdot y=y.
\]
Así, $r\cdot y=y$ para toda cara $y\in\mathcal F$.

Esto significa que $r$ fija cada cara del cubo. Pero la única rotación que fija todas las caras es la identidad, contradiciendo que $r$ es no trivial. Luego
\[
X^r=\varnothing
\qquad\text{para toda }r\in R\setminus\{\operatorname{id}\}.
\]

\section*{5. Aplicación del lema de Burnside}
Por el lema de Burnside, el número de órbitas es
\[
\frac{1}{|R|}\sum_{r\in R}|X^r|
=\frac{1}{24}\left(6!+23\cdot0\right)
=\frac{6!}{24}
=\boxed{30}.
\]
\hfill$\square$
\end{document}
